% F16 fig:newton (Chapter 7, section 7.4): Newton's method for T with H(T) = H(1000 K) from 800 K, the slope by
% Central<2> with the default step (E12). The residual |H(T) - H*| in J/mol against the iteration: the number of
% correct digits doubles per step until the residual reaches the rounding of H. Data: DS11, data/newton.csv.
\begin{tikzpicture}
\begin{axis}[eqtufte, width=38mm, height=34mm, ymode=log,
  xlabel={iteration}, ylabel={residual $|H(T) - H^\ast|$ (J/mol)},
  xtick distance=1, ytick={1e-12,1e-8,1e-4,1,1e4},
  ymin=1e-13, ymax=1e5, xmin=-0.2, xmax=4.4]
  \addplot[accent, mark=*, mark size=1pt]
    table[col sep=comma, x=iteration, y expr={abs(\thisrow{residual})}, comment chars={\#}] {data/newton.csv}
    node[pos=0.5, anchor=south west, font=\scriptsize, accent, align=left] {quadratic:\protect\\digits double}
    node[pos=1, left=1pt, font=\scriptsize, muted] {rounding of $H$};
\end{axis}
\end{tikzpicture}
